\documentclass[aps,prb,onecolumn,nofootinbib,superscriptaddress,floatfix]{revtex4-2}
\usepackage{amsmath,amssymb,amsthm,mathtools,bm}
\usepackage{booktabs}
\usepackage[colorlinks=true,linkcolor=blue,citecolor=blue,urlcolor=blue]{hyperref}
\usepackage{microtype}
\hypersetup{pdftitle={Channel relaxation and stationary-state occupation gaps},
  pdfauthor={Asadullah Bhuiyan}}
\newcommand{\Id}{\hat{\mathbf 1}}
\newcommand{\id}{\mathbf 1}
\newcommand{\Tr}{\operatorname{Tr}}
\newcommand{\tr}{\operatorname{tr}}
\newcommand{\E}{\widetilde{\mathcal E}}
\newcommand{\D}{\mathcal D}
\newcommand{\spec}{\operatorname{spec}}
\newcommand{\steady}{\mathrm{ss}}
\begin{document}
\title{Channel relaxation and stationary-state occupation gaps}
\author{Asadullah Bhuiyan}
\noaffiliation
\date{October 6, 2026}
\begin{abstract}
We distinguish the relaxation spectrum of an outcome-averaged fermionic
channel from the occupation spectrum of its stationary correlation
matrix. For a fixed sequence of perfect occupation resets, we derive
the complete charge-neutral correlation hierarchy and show that the
neutral and parity-even many-body channel gaps are exactly
$1-\rho(A)^2$, where $A$ is the ordered product of complementary
single-particle projectors. This equality requires more than closure
of the two-point equation. Explicit local and nonlocal examples
separate stationary purity gaps from dynamical gaps. We also establish
the gap relations under probabilistic application, Poissonization, and
finite-time Lindblad evolution, with their sector and locality
qualifications. Verified hard-wall data illustrate the distinction
for fixed-width and square geometries; finite-size occupation gaps
are not interpreted as thermodynamic extrapolations.
\end{abstract}
\maketitle
\setcounter{tocdepth}{0}
\makeatletter
\let\l@subsection\@gobbletwo
\makeatother
\tableofcontents
\appendix

\section{Elementary resets and covariance evolution}
\label{sec:reset}
There are two different questions behind the word \emph{gap}. One asks
how rapidly a perturbation of a state decays under repeated evolution.
The other asks whether the stationary correlation matrix has an
occupation close to $1/2$. The first concerns the map that evolves a
state; the second concerns the state reached by that map. We will
derive both objects explicitly rather than infer one spectrum from
the other.

The argument proceeds in three stages. First, we derive the action
of one measurement-and-feedback step and compose the steps into a
cycle. Second, we identify the covariance relaxation spectrum and
prove when it determines the many-body gap. Third, we compare this
dynamical information with stationary occupations and explain what
changes on passing to continuous time. Simple matrix examples will
make the distinctions concrete before we return to the numerical data.

We consider the physical-layer channel obtained after discarding the
measurement record and fresh ancillary modes of the adaptive circuit
of Bhuiyan, Pan, and Jian (BPJ)~\cite{BPJ}. The assumptions throughout
Apps.~\ref{sec:reset} and \ref{sec:spectrum} are perfect feedback,
normalized measured modes, and a fixed sequence repeated each cycle.
The measured overcomplete Wannier (OW) modes need not be orthogonal.
Support truncation is allowed, provided each retained mode is
normalized. We retain the measurement dephasing already present in
the reset channel; we do not append an additional dephasing operation.

\subsection{Kraus representation and exact dissipator identity}
Let $\{\hat\psi_a,\hat\psi_b^\dagger\}=\delta_{ab}$ for $N$ fermion
modes, and define
\begin{equation}
 \hat\chi_j=\sum_a v_{ja}\hat\psi_a,\qquad
 v_j^\dagger v_j=1,\qquad \hat N_j=\hat\chi_j^\dagger\hat\chi_j.
 \label{eq:mode}
\end{equation}
Here $v_j$ is a column of coefficients: the definition of
$\hat\chi_j$ contains no implicit complex conjugation. Emptying and
filling this mode correspond, respectively, to
\begin{align}
 \E_{j,0}(\hat\rho)&=\hat\chi_j\hat\rho\hat\chi_j^\dagger
 +(\Id-\hat N_j)\hat\rho(\Id-\hat N_j),\nonumber\\
 \E_{j,1}(\hat\rho)&=\hat\chi_j^\dagger\hat\rho\hat\chi_j
 +\hat N_j\hat\rho\hat N_j. \label{eq:kraus}
\end{align}
Since $\hat N_j^2=\hat N_j$ and
$\hat\chi_j\hat\chi_j^\dagger=\Id-\hat N_j$, each pair of Kraus
operators satisfies $\sum_m\hat K_m^\dagger\hat K_m=\Id$.
These are completely positive trace-preserving maps.

The meaning of a reset can be seen without any lattice structure.
For a single mode, in the empty/occupied basis, write
\begin{equation}
 \hat\rho=
 \begin{pmatrix}1-f&z\\z^*&f\end{pmatrix},
 \qquad
 \E_0(\hat\rho)=
 \begin{pmatrix}1&0\\0&0\end{pmatrix},
 \qquad
 \E_1(\hat\rho)=
 \begin{pmatrix}0&0\\0&1\end{pmatrix}.
 \label{eq:one_mode_reset}
\end{equation}
Here $f$ is the initial occupation and $z$ is a coherence; their
allowed values are fixed by positivity of $\hat\rho$. The reset
forgets both. A physical parity-even single-mode state has $z=0$;
retaining $z$ here also displays the action on the formal odd sector.
In the emptying map, an occupied input loses its
particle through $\hat\chi$, while an empty input is left empty by
$\Id-\hat N$. The filling map exchanges the two roles. Summing
the branches discards the record, so the two alternatives are added
as density matrices, not as coherent amplitudes.

With the normalization
\begin{equation}
 \D[\hat L](\hat\rho)=\hat L\hat\rho\hat L^\dagger
 -\tfrac12\{\hat L^\dagger\hat L,\hat\rho\},
 \label{eq:dissipator}
\end{equation}
expanding the projector terms in Eq.~\eqref{eq:kraus} gives
\begin{align}
 \E_{j,0}&=\mathrm{Id}+\D[\hat\chi_j]+\D[\hat N_j],\nonumber\\
 \E_{j,1}&=\mathrm{Id}+\D[\hat\chi_j^\dagger]+\D[\hat N_j].
 \label{eq:idplusL}
\end{align}
For example, the two anticommutators in the emptying case sum to
$-\{\hat N_j,\hat\rho\}$, while the remaining terms are
$\hat\chi_j\hat\rho\hat\chi_j^\dagger+\hat N_j\hat\rho\hat N_j$.
Adding $\hat\rho$ reproduces the first line of Eq.~\eqref{eq:kraus}.
The filling case follows using $\hat\chi_j\hat\chi_j^\dagger=\Id-\hat N_j$.
Equation~\eqref{eq:idplusL} is an exact discrete-channel identity,
not a short-time approximation or a claim that $\E_j=e^{\mathcal L_j}$.

\subsection{Closed two-point equation}
For the outcome-averaged state define
\begin{equation}
 \overline G_{ab}=\Tr(\overline{\hat\rho}\,
             \hat\psi_a^\dagger\hat\psi_b),\qquad
 P_j=v_jv_j^\dagger,\quad Q_j=\id-P_j. \label{eq:G}
\end{equation}
The total charge is $\langle\hat N_F\rangle=\tr\overline G$, with
$\tr$ the single-particle trace and $\Tr$ the many-body trace.
No Gaussian-state assumption enters this definition. The canonical
anticommutation relations give
\begin{equation}
 [\hat N_j,\hat\psi_b]=-v_{jb}^*\hat\chi_j,
 \qquad [\hat N_j,\hat\psi_a^\dagger]=v_{ja}\hat\chi_j^\dagger.
 \label{eq:commutators}
\end{equation}
Using these identities in
$\D^\dagger[\hat N_j](\hat O)
=-\tfrac12[\hat N_j,[\hat N_j,\hat O]]$, and evaluating the
gain/loss term on $\hat O=\hat\psi_a^\dagger\hat\psi_b$,
gives the two contributions
\begin{align}
 (\delta\overline G)_{\rm gain/loss}
   &=-\tfrac12(P_j\overline G+\overline GP_j)+s_jP_j,\nonumber\\
 (\delta\overline G)_{\rm deph}
   &=P_j\overline GP_j-\tfrac12(P_j\overline G+\overline GP_j).
 \label{eq:two_contributions}
\end{align}
Thus the exact elementary update is
\begin{equation}
 \overline G'=Q_j\overline GQ_j+s_jP_j,
 \qquad s_j\in\{0,1\}. \label{eq:resetG}
\end{equation}
In a basis beginning with the measured mode, this erases its cross
correlations and sets its occupation to $s_j$, leaving the
orthogonal block unchanged. Explicitly, in that basis,
\begin{equation}
 \overline G=
 \begin{pmatrix} f&b^\dagger\\ b&G_\perp\end{pmatrix}
 \quad\longmapsto\quad
 \overline G'=
 \begin{pmatrix} s_j&0\\0&G_\perp\end{pmatrix}.
 \label{eq:block_reset}
\end{equation}
The complementary projector $Q_j$ therefore retains only the part
of the old correlations orthogonal to the measured mode, while
$s_jP_j$ inserts the new occupation. Equations~\eqref{eq:mode} and
\eqref{eq:G} fix the conjugation convention needed to use $P_j=v_jv_j^\dagger$.

One cycle is the ordered composition
\begin{equation}
 \E=\prod_{j=M}^{1}\E_{j,s_j}
   \equiv\E_{M,s_M}\circ\cdots\circ\E_{1,s_1},
 \label{eq:product}
\end{equation}
where step 1 acts first. Iteration of Eq.~\eqref{eq:resetG} yields
\begin{align}
 \overline G'&=A\overline GA^\dagger+B,\nonumber\\
 A&=Q_M\cdots Q_1,\qquad
 B=\sum_{j=1}^{M}s_jR_jP_jR_j^\dagger,\nonumber\\
 R_j&=Q_M\cdots Q_{j+1},\qquad R_M=\id.
 \label{eq:affine}
\end{align}
Noncommuting projectors make the order physical. The targets $s_j$
enter $B$, whereas $A$ depends only on the measured modes and order.
The word \emph{affine} refers to the state-independent source $B$:
the update is a linear transformation plus an additive term.
After $n$ complete cycles, it reads
\begin{equation}
 \overline G_n=A^n\overline G_0(A^\dagger)^n+
       \sum_{\ell=0}^{n-1}A^\ell B(A^\dagger)^\ell .
 \label{eq:finite_cycles}
\end{equation}
The first term remembers the initial state; the sum collects
occupations injected during earlier cycles and subsequently
propagated by the remaining resets.
In particular, the stationary equation is affine:
\begin{equation}
 G_\steady=AG_\steady A^\dagger+B. \label{eq:fixedG}
\end{equation}
For $r\equiv\rho(A)<1$ its unique solution is
$G_\steady=\sum_{n=0}^{\infty}A^nB(A^\dagger)^n$.
Stationarity therefore does not require an eigenvalue 1 of $A$.
This point is essential: the many-body channel has a stationary
eigenvector, but the smaller matrix $A$ describes the survival of
initial correlations, not the entire state including its source.

\section{Relaxation spectrum and many-body gap}
\label{sec:spectrum}
\subsection{Definitions and spectral inclusion}
An eigenoperator $\hat X$ of the channel satisfies
$\E(\hat X)=\mu\hat X$. After $n$ cycles its amplitude is multiplied
by $\mu^n$. Thus $|\mu|$ determines decay and the phase of $\mu$
determines oscillation. A multiplier $0.9$ decays more slowly than
$0.2$; a multiplier $-0.9$ has the same decay envelope as $0.9$
but alternates in sign. The slowest relaxing mode is therefore the
one with modulus closest to 1, not the smallest eigenvalue.

Fix an invariant linear operator sector $\mathcal V$ containing a
stationary density matrix. The channel preserves the trace-zero
subspace $\mathcal V_0$. This is the natural space of perturbations:
the difference of any two normalized density matrices has zero trace.
In particular, writing $\hat\rho=\hat\rho_\steady+\delta\hat\rho$ gives
\begin{equation}
 \E^n(\hat\rho)=\hat\rho_\steady+\E^n(\delta\hat\rho),
 \qquad \Tr(\delta\hat\rho)=0.
 \label{eq:state_perturbation}
\end{equation}
The first term is supposed to remain; only the second must decay
for relaxation to occur. Excluding the stationary eigenvalue does
not hide a slow perturbation: it removes the reference state itself.
Define the absolute relaxation gap and
logarithmic rate by
\begin{align}
 r_{\mathcal V}&=\rho(\E|_{\mathcal V_0}),&
 g_{\mathcal V}&=1-r_{\mathcal V},\nonumber\\
 \kappa_{\mathcal V}&=-\log r_{\mathcal V}.&&
 \label{eq:gapdef}
\end{align}
When $r_{\mathcal V}<1$, the stationary eigenvalue is simple in this
sector and all other eigenvalues lie strictly inside the unit circle;
the channel is mixing~\cite{Burgarth,Wolf}. Removing the stationary
mode means restricting away from this one fixed direction, not
discarding all troublesome eigenvalues. An additional stationary
state supplies a trace-zero fixed direction, and any other
unit-modulus eigenvalue also gives $g_{\mathcal V}=0$. One can instead
define a gap above an entire fixed-point space, but that is not a
claim of relaxation to a unique state.

The two gap conventions encode the same multiplier in different
units. If $r_{\mathcal V}=1-g_{\mathcal V}$, the slow envelope is
$r_{\mathcal V}^n=e^{-n\kappa_{\mathcal V}}$, and
\begin{equation}
 \kappa_{\mathcal V}=-\log(1-g_{\mathcal V})
     =g_{\mathcal V}+\tfrac12g_{\mathcal V}^2+\cdots.
 \label{eq:log_conversion}
\end{equation}
The dimensionless channel gap and the decay rate per cycle agree
only near a closing. For a large gap, replacing one by the other
can substantially change the numerical value.

The absolute gap is $1-\max|\mu|$, not $\min|1-\mu|$.
For a Lindbladian $\mathcal L$, the corresponding decay rate is
\begin{equation}
 \delta_{\mathcal L,\mathcal V}
   =-\max_{\lambda\in\spec(\mathcal L|_{\mathcal V_0})}
        \operatorname{Re}\lambda. \label{eq:Lgap}
\end{equation}
The stationary zero is excluded under the same sector convention.
Additional zeros or purely imaginary modes obstruct mixing.

There are three matrix sizes to distinguish. For $N$ fermion modes,
$A$ and $\overline G$ are $N\times N$ matrices, a general linear map
on their entries is $N^2\times N^2$, and the unrestricted many-body
channel acts on an operator space of dimension $4^N$.
The density matrix itself is $2^N\times2^N$. The simplification
is that a tractable part of the exponentially larger evolution
closes on the two-point function.

Subtracting Eq.~\eqref{eq:fixedG} gives
\begin{equation}
 \delta G'=A\delta GA^\dagger,\qquad
 \operatorname{vec}(\delta G')=(A^*\otimes A)
                                      \operatorname{vec}(\delta G).
 \label{eq:vec}
\end{equation}
Column vectorization means stacking the columns of $\delta G$ into
one vector. It is a bookkeeping device, not another physical
approximation. To see the spectrum directly, suppose first that
$Aw_i=a_iw_i$ gives an eigenbasis. Then
\begin{equation}
 A(w_iw_j^\dagger)A^\dagger
     =a_i a_j^*\,w_iw_j^\dagger.
 \label{eq:pairwise_action}
\end{equation}
Each side of the covariance contributes one factor. Schur
triangularization gives the same eigenvalues even without a
complete eigenbasis. The covariance multipliers are therefore
$a_i a_j^*$, including the positive multiplier $|a_i|^2$.
Consequently
\begin{equation}
 g_C=1-r^2,\qquad \kappa_C=-2\log r,\qquad r=\max_i|a_i|.
 \label{eq:covgap}
\end{equation}
No eigenvalue of $A$ is removed in computing $r$. The square comes
from the two-sided action, not from a choice of gap convention.
The spectral radius, rather than the largest singular value, controls
asymptotic decay. Below we take $0<r<1$; if $r=0$, the finite-dimensional
homogeneous map is nilpotent and its logarithmic rate is infinite.

\subsection{A two-mode example of the projector product}
For a concrete illustration, empty two real normalized modes in
sequence, with coefficient vectors
$v_1=(1,0)^{\mathsf T}$ and
$v_2=(\cos\theta,\sin\theta)^{\mathsf T}$.
Writing $c_\theta=\cos\theta$ and $s_\theta=\sin\theta$, their
complementary projectors and ordered product are
\begin{equation}
 Q_1=\begin{pmatrix}0&0\\0&1\end{pmatrix},\qquad
 Q_2=\begin{pmatrix}s_\theta^2&-c_\theta s_\theta\\
              -c_\theta s_\theta&c_\theta^2\end{pmatrix},\qquad
 A=Q_2Q_1=\begin{pmatrix}0&-c_\theta s_\theta\\
                          0&c_\theta^2\end{pmatrix}.
 \label{eq:two_mode_product}
\end{equation}
The eigenvalues of $A$ are $0,c_\theta^2$, whereas the covariance
map has eigenvalues $0,0,0,c_\theta^4$. Hence
$r=c_\theta^2$ and $g_C=1-c_\theta^4$. Nearly parallel measured
modes leave an almost unmeasured direction and give slow
relaxation. Orthogonal measured modes erase all initial
correlations in one cycle. At $\theta=0$ the second physical mode
is never measured, and the gap vanishes.

Both targets in this example are empty, so $B=0$. For
$0<\theta<\pi/2$ the stationary matrix is $G_\steady=0$ even
though the relaxation rate varies with $\theta$. The example also
separates eigenvalues from singular values: the largest singular
value of $A$ is $|c_\theta|$, not $c_\theta^2$. A largest-singular-value
calculation would therefore not give the same asymptotic rate.

\subsection{Why covariance closure initially gives only a bound}
Covariance closure alone proves an upper bound. The adjoint channel
preserves the span of $\Id$ and the bilinears
$\hat\psi_a^\dagger\hat\psi_b$. After centering about $G_\steady$,
the covariance eigenvalues are contained in the many-body spectrum.
Thus, for any mixing sector containing these operators,
\begin{equation}
 r_{\mathcal V}\ge r^2,\qquad
 g_{\mathcal V}\le g_C,\qquad
 \kappa_{\mathcal V}\le\kappa_C. \label{eq:bound}
\end{equation}
A positive upper bound cannot rule out a vanishing many-body gap.
Knowing one family of decaying observables need not identify the
slowest observable in the entire system. An unseen higher-order
correlation could decay more slowly. To replace the inequality
by an equality, we must rule out precisely that possibility.
The continuous-time correlation hierarchy of Barthel and
Zhang~\cite{BZ}, especially their Proposition 7, provides the analogous
spectral-inclusion reasoning. Equality here will follow from the
specific reset structure, not from closure by itself.

The covariance mode can also be exhibited as an observable.
For $A^\dagger u=a^*u$, define
\begin{equation}
 \hat O_u=\sum_{ab}u_a^*u_b
   \left[\hat\psi_a^\dagger\hat\psi_b-(G_\steady)_{ab}\Id\right].
 \label{eq:observable}
\end{equation}
Since its expectation is $u^\dagger\delta Gu$,
$\E^\dagger(\hat O_u)=|a|^2\hat O_u$.
Choosing $|a|=r$ realizes the slow neutral observable. A right
eigenvector $Aw=aw$ instead gives the right covariance eigenmatrix
$ww^\dagger$; these left and right constructions should not be
interchanged for nonnormal $A$.

\subsection{Higher moments and exterior powers}
Within the parity-even sector, a perfect reset never creates
a higher-degree observable from a lower-degree one. We use
\emph{degree} to count creation and annihilation operators in a
normal-ordered monomial. Bilinears have degree two; the corresponding
four-point moments have degree four. We keep the full moments,
not a Gaussian factorization of them.

Choose a single-particle basis beginning with $\hat\chi_j$, and
let $\hat B$ contain only orthogonal modes. On even total operators,
direct substitution into Eq.~\eqref{eq:kraus} gives
\begin{align}
 \E_{j,s}^\dagger(\hat B)&=\hat B
       &&(\hat B\text{ even}),\nonumber\\
 \E_{j,s}^\dagger(\hat\chi_j^{(\dagger)}\hat B)&=0
       &&(\hat B\text{ odd}),\nonumber\\
 \E_{j,s}^\dagger(\hat N_j\hat B)&=s\hat B
       &&(\hat B\text{ even}). \label{eq:localrule}
\end{align}
A single measured leg is erased, while a measured occupation pair
is replaced by its target and lowers the degree by two. This is an
operator identity and makes no use of Wick factorization.

For increasing $p$-element index sets $I,J$, introduce
\begin{equation}
 F_{I;J}^{(p)}=\Tr\!\left(\hat\rho\,
 \hat\psi_{i_1}^\dagger\cdots\hat\psi_{i_p}^\dagger
 \hat\psi_{j_p}\cdots\hat\psi_{j_1}\right).
 \label{eq:moments}
\end{equation}
For $p=1$ these are the entries of $\overline G$; for $p=2$ they are
four-point moments. In vector form, the first levels of their
update have the schematic structure
\begin{equation}
 \begin{pmatrix}1\\F^{(1)}\\F^{(2)}\\\vdots\end{pmatrix}'=
 \begin{pmatrix}
  1&0&0&\cdots\\
  b_1&H_1&0&\cdots\\
  b_2&T_{21}&H_2&\cdots\\
  \vdots&\vdots&\vdots&\ddots
 \end{pmatrix}
 \begin{pmatrix}1\\F^{(1)}\\F^{(2)}\\\vdots\end{pmatrix}.
 \label{eq:triangular_example}
\end{equation}
Here $b_p$ and $T_{21}$ denote lower-order source terms; their
values are not needed to determine the eigenvalues of this
block-triangular matrix. What matters is the diagonal block $H_p$
propagating the degree-$2p$ part.

The homogeneous degree-$2p$ term is obtained by projecting each leg
onto the orthogonal complement of the measured mode.
Antisymmetrization turns the $p$ creation-leg factors into the
exterior power $\bigwedge^p Q_j$ and the annihilation-leg factors
into its conjugate. The one-cycle homogeneous block is therefore
\begin{equation}
 H_p[I,J;K,L]=\det A[I,K]\,\det A[J,L]^*.
 \label{eq:minors}
\end{equation}
The exterior-power notation is simply the antisymmetrized action
on several fermionic indices. For example, on two creation legs,
\begin{equation}
 \left(\bigwedge\nolimits^2 A\right)_{\{i,j\},\{k,\ell\}}
       =A_{ik}A_{j\ell}-A_{i\ell}A_{jk}.
 \label{eq:wedge_example}
\end{equation}
The minus sign is the fermionic exchange sign. For $p$ legs,
the same sum over signed permutations is the $p\times p$
determinant in Eq.~\eqref{eq:minors}.
Indeed, the Cauchy--Binet identity composes the minors of successive
$Q_j$ in exactly the order in Eq.~\eqref{eq:affine}. All source
terms from Eq.~\eqref{eq:localrule} belong to lower-degree blocks.
The full hierarchy, ordered by $p$, is block triangular.

The eigenvalues of $\bigwedge^p A$ are the products of $p$ distinct
eigenvalues of $A$. This remains true for defective $A$, as follows
by taking exterior powers of its Schur-triangular form. Hence the
neutral-sector spectrum, with algebraic multiplicities, is
\begin{equation}
 \Lambda_{I,J}=\prod_{i\in I}a_i\prod_{j\in J}a_j^*,\qquad
 |I|=|J|=p,\quad 0\le p\le N. \label{eq:hierarchy}
\end{equation}
The sets $I$ and $J$ are independently chosen and may overlap.
In particular, the same eigenvalue may appear on both legs.

To see that no neutral operators are omitted, the normal-ordered
monomials with equal numbers of creation and annihilation operators
form a basis. One may construct each number-preserving Fock-space
matrix unit by multiplying its creation and annihilation string by
occupation or vacancy projectors on the remaining sites. Expanding
these projectors gives precisely such monomials. Their number is
\begin{equation}
 \sum_{p=0}^{N}\binom Np^2=\binom{2N}{N}, \label{eq:span}
\end{equation}
which is also the dimension of all operators block diagonal in
particle number. Spanning and this dimension count establish
completeness of Eq.~\eqref{eq:hierarchy}.

The constant block $p=0$ contributes one eigenvalue 1. For $p\ge1$,
$|\Lambda_{I,J}|\le r^{2p}\le r^2$, and the bound is attained at
$p=1$ with $I=J$ corresponding to $|a_i|=r$. Thus the stationary
eigenvalue is simple and
\begin{equation}
 \boxed{g_{q=0}=g_C=1-r^2,\qquad \kappa_{q=0}=\kappa_C.}
 \label{eq:neutral}
\end{equation}
Higher neutral correlations cannot have a slower spectral decay.

\subsection{Neutral, even, and odd sectors}
Charge neutrality means $[\hat N_F,\hat\rho]=0$, with
$\hat N_F=\sum_a\hat\psi_a^\dagger\hat\psi_a$.
It permits mixtures of different particle numbers. The reset channel
preserves neutrality, but it does not conserve particle number:
its Kraus operators have definite, not necessarily zero, charge.
Both maximally mixed initial states and fixed-number Slater states
belong to the neutral sector.

Examples clarify the distinction between charge and parity.
$\hat\psi_1^\dagger\hat\psi_2$ is neutral because it has one
creation and one annihilation operator.
$\hat\psi_1\hat\psi_2$ is not neutral, but it is parity even
because it contains two fermionic operators.
$\hat\psi_1$ is parity odd.
An invariant operator sector groups quantities whose evolution
stays within that group; it is not necessarily a fixed-particle-number
subspace of state vectors.

For a monomial with $p$ creation and $\ell$ annihilation operators,
the same even-operator argument gives a homogeneous block
$\bigwedge^p A\otimes(\bigwedge^\ell A)^*$, up to the chosen
vectorization order. Every nonconstant even monomial has
$p+\ell\ge2$. Its eigenvalue modulus is at most $r^{p+\ell}\le r^2$,
and the neutral bilinears attain $r^2$. Consequently
\begin{equation}
 g_{\rm even}=g_{q=0}=g_C,\qquad
 \kappa_{\rm even}=\kappa_C. \label{eq:even}
\end{equation}

For completeness, the specified Kraus maps also define an extension
to parity-odd operators. Let $\hat P_F=(-1)^{\hat N_F}$ and
$\mathcal S(\hat O)=\hat P_F\hat O$, so $\mathcal S^2=\mathrm{Id}$.
Writing $\hat J=\hat\chi$ or $\hat\chi^\dagger$ for the odd Kraus
operator and $\hat R$ for the even projector Kraus operator,
\begin{equation}
 (\mathcal S\E_j^\dagger\mathcal S)(\hat O)
   =-\hat J^\dagger\hat O\hat J+\hat R\hat O\hat R.
 \label{eq:odd_similarity}
\end{equation}
On odd operators, this extra minus sign cancels the sign from moving
an odd orthogonal-mode string past $\hat J$. The transformed map
therefore leaves an orthogonal odd string unchanged, erases a
single measured leg, and sends $\hat N\hat B$ to $s\hat B$.
Its homogeneous hierarchy has odd $p+\ell$, with maximal modulus
$r$ attained by one-leg blocks. Similarity preserves the spectrum
and respects composition across the cycle. On unrestricted
operator space,
\begin{equation}
 g_{\rm full}=1-r,\qquad
 \kappa_{\rm full}=-\log r=\kappa_C/2. \label{eq:full}
\end{equation}
Physical fermionic density matrices are parity even; Eq.~\eqref{eq:even}
is the relevant many-body gap for the preparations considered here.
Equation~\eqref{eq:full} explains why an unrestricted spectral
definition can differ without contradicting covariance closure.

These equalities are specific to the exact reset hierarchy.
Imperfect feedback, extra dephasing, or averaging over a new random
schedule each cycle requires reexamining the higher moments.
Nonnormality can also produce long transients and Jordan-polynomial
prefactors. A spectral gap fixes the asymptotic exponential envelope
at fixed size, not a size-uniform mixing time or a conditioned
trajectory-purification time.

\section{Stationary occupation gaps and counterexamples}
\label{sec:stationary}
\subsection{Distinct spectra}
Let $n_a\in[0,1]$ be the eigenvalues of the full-system stationary
correlation matrix $G_\steady$. We use
\begin{equation}
 \Delta=\min_a|n_a-\tfrac12|,\qquad
 \Delta_{\rm pur}=2\Delta=\min_a|1-2n_a|.
 \label{eq:purity}
\end{equation}
These are stationary occupation diagnostics; they are not eigenvalue
gaps of a channel. Our normalization makes the purity gap unity for
a pure number-conserving Gaussian state and zero when an occupation
equals $1/2$. Other normalizations, including squared conventions,
must be converted before comparison. The full-system matrix is also
distinct from a subsystem restriction used for an entanglement spectrum.

The reference point $1/2$ comes from the occupation of a maximally
mixed fermion, not from a stationary eigenvalue of the evolution.
For example, occupations $0.01$ and $0.99$ are far from maximal
mixing and give $\Delta_{\rm pur}=0.98$, while an occupation
$0.499$ gives $\Delta_{\rm pur}=0.002$. Neither statement says how
long it took to prepare those occupations. Also, the symbol
$\rho(A)$ denotes a \emph{spectral radius}; it is not a density
matrix $\hat\rho$. Likewise, $n_a$ labels occupations of $G_\steady$,
whereas $\mu$ labels multipliers of a dynamical map.

For a Gaussian state the single-particle modular energies are
\begin{equation}
 \epsilon_a=\log\frac{1-n_a}{n_a}. \label{eq:modular}
\end{equation}
For a non-Gaussian outcome-averaged state, this defines the Gaussian
reference state sharing its two-point function, not the spectrum of
its actual many-body modular Hamiltonian $-\log\hat\rho_\steady$.
An occupation at $1/2$ denotes maximal mixing of that covariance
mode; it does not identify a conserved observable or a dark mode
of the evolution.

The separation is already visible in Eq.~\eqref{eq:affine}:
relaxation of $\delta G$ depends on $A$, while $G_\steady$ also depends
on $B$. For gauge-invariant quasi-free channels, complete positivity
in this convention requires $0\le B\le\id-AA^\dagger$~\cite{DFP}.
For example, $A=\sqrt q\,\id$ and $B=(1-q)G_0$, with
$0\le G_0\le\id$ and $0\le q<1$, give $G_\steady=G_0$ and
$g_C=1-q$, independently of the occupation gap of $G_0$.
For a single occupation, the same construction is the elementary
recursion
\begin{equation}
 n_{m+1}=q n_m+(1-q)n_*,
 \qquad n_m-n_*=q^m(n_0-n_*).
 \label{eq:scalar_source}
\end{equation}
The number $n_*$ specifies where the evolution ends; $q$ specifies
how much of the deviation survives each step. One may change
$n_*$ from nearly empty to exactly half-filled without changing
$q$. This is the simplest way to see why a purity-gap closing does
not force slow relaxation.

This Gaussian-channel construction is illustrative; closure of the
two-point equation does not itself make our full reset channel
Gaussian on arbitrary inputs.

\subsection{Global and local mixing examples}
For any density matrix $\hat\sigma$, let
$\mathcal T_\sigma(\hat X)=\Tr(\hat X)\hat\sigma$.
The replacement channel
\begin{equation}
 \E_p=(1-p)\mathrm{Id}+p\mathcal T_\sigma,\qquad 0<p\le1,
 \label{eq:replace}
\end{equation}
has the unique stationary state $\hat\sigma$ and multiplier $1-p$
on every trace-zero operator. Its gap is $p$, independently of the
stationary spectrum. To verify the claim, write
$\hat\rho=\hat\sigma+\delta\hat\rho$ with
$\Tr\delta\hat\rho=0$. Since
$\mathcal T_\sigma(\delta\hat\rho)=0$, the channel maps the
perturbation to $(1-p)\delta\hat\rho$ while leaving $\hat\sigma$
unchanged. There is no additional slow subspace to inspect.
Taking $\hat\sigma=\Id/d$ gives a maximally
mixed endpoint and, for a fermionic Fock space, $G_\steady=\id/2$.
The generator $\gamma(\mathcal T_\sigma-\mathrm{Id})$ has nonstationary
eigenvalues $-\gamma$, and
\begin{equation}
 e^{t\gamma(\mathcal T_\sigma-\mathrm{Id})}
   =\mathcal T_\sigma+e^{-\gamma t}
           (\mathrm{Id}-\mathcal T_\sigma). \label{eq:replace_semigroup}
\end{equation}
Global replacement is generally nonlocal, so it alone does not
address locality constraints.

A strictly local example is independent depolarization of spins.
Let $\mathcal T_i$ replace spin $i$ by $\id_2/2$ while tracing out
its old state. Then
\begin{equation}
 \mathcal L_{\rm loc}=\gamma\sum_i(\mathcal T_i-\mathrm{Id})
   =\frac{\gamma}{4}\sum_{i,a=x,y,z}\D[\hat\sigma_i^a].
 \label{eq:local_depol}
\end{equation}
A Pauli string of weight $w$ has eigenvalue $-\gamma w$.
The weight counts sites with a nonidentity Pauli factor.
Replacing any such site erases that string, so each of its $w$
sites contributes one decay rate $\gamma$. The identity string
has weight zero and is stationary; the slowest nonconstant strings
have weight one.
The unique stationary state is maximally mixed and the gap is
$\gamma$, uniformly in the number of spins. Sampling this local
semigroup at time $\tau$ gives a channel gap $1-e^{-\gamma\tau}$.
No locality-preserving spin-to-fermion mapping is needed or asserted.

For a single fermionic mode, balanced gain and loss give
\begin{equation}
 \mathcal L_{\rm bal}=\eta\bigl(\D[\hat c]+\D[\hat c^\dagger]\bigr),
 \qquad \dot n=\eta(1-2n). \label{eq:balanced}
\end{equation}
Its full spectrum is $\{0,-\eta,-\eta,-2\eta\}$ and its stationary
state is $\Id/2$. The even population sector decays at $2\eta$;
the unrestricted gap is $\eta$. Thus even the simplest fermionic
example requires stating which operator sector is meant.

\subsection{Period-two channel and stationary-mode exclusion}
The channel
\begin{equation}
 \E_{\rm flip}(\hat\rho)=\hat c\hat\rho\hat c^\dagger
                         +\hat c^\dagger\hat\rho\hat c
 \label{eq:flip}
\end{equation}
is trace preserving because $\hat c^\dagger\hat c+
\hat c\hat c^\dagger=\Id$. Its action on an arbitrary matrix is
\begin{equation}
 \E_{\rm flip}\!
 \begin{pmatrix}x&z\\w&y\end{pmatrix}
       =\begin{pmatrix}y&0\\0&x\end{pmatrix}.
 \label{eq:flip_matrix}
\end{equation}
The diagonal entries are exchanged, while off-diagonal entries are
erased. A fixed matrix must have $x=y$ and $z=w=0$, proving
directly that the stationary space is spanned only by $\Id$.
On $\Id,\hat\sigma_z,
\hat\sigma_x,\hat\sigma_y$ it has eigenvalues
$1,-1,0,0$, respectively. Its stationary space is one dimensional,
spanned by $\Id$, yet it is not mixing: a population imbalance changes
sign every step. The absolute channel gap is zero because of $-1$.
The two zero multipliers instead describe coherences erased in one
step. Neither taking the smallest eigenvalue nor simply deleting
all unit-modulus eigenvalues gives the relaxation gap.

Its Poissonized generator $\gamma(\E_{\rm flip}-\mathrm{Id})$
has spectrum $\{0,-2\gamma,-\gamma,-\gamma\}$ and is precisely
Eq.~\eqref{eq:balanced} with $\eta=\gamma$. Randomizing the number
of flips removes the period-two obstruction without changing the
stationary state. This supplies a counterexample to a general
reverse implication from a gapped Poissonized generator to a
mixing discrete channel.

Finally, a gapped Hamiltonian can have a maximally mixed
infinite-temperature Gibbs state: for example
$\hat H=\varepsilon\sum_i\hat n_i$ has excitation gap $\varepsilon$
but $\hat\rho_{\beta=0}=\Id/2^N$. This only distinguishes a
Hamiltonian excitation gap from a state's occupation gap; unitary
evolution under that Hamiltonian does not prepare this state as a
unique attractor.

\section{Continuous-time connection}
\label{sec:continuous}
\subsection{Probabilistic application and Poissonization}
For an arbitrary channel $\E$, define
\begin{equation}
 \E_p=(1-p)\mathrm{Id}+p\E,\qquad 0<p\le1. \label{eq:lazy}
\end{equation}
Its fixed-point space is exactly that of $\E$, since
$\E_p-\mathrm{Id}=p(\E-\mathrm{Id})$. Each eigenvalue transforms
as $\mu_p=1-p+p\mu$, with no diagonalizability assumption.
If $g=1-\max_{\mathcal V_0}|\mu|>0$, then
\begin{equation}
 |\mu_p|\le1-p+p|\mu|\le1-pg,\qquad
 \boxed{g(p)\ge pg>0.} \label{eq:lazy_bound}
\end{equation}
Hence the mixing gap cannot close at any fixed $p>0$ along this
particular deformation. The identity endpoint $p=0$ is excluded:
every state is then stationary. This is a spectral estimate, not
an argument from continuity alone.

Taking $p=\gamma\,dt$ gives the generator
\begin{equation}
 \mathcal L_{\rm cyc}=\gamma(\E-\mathrm{Id})
   =\sum_a\D[\sqrt\gamma\,\hat K_a], \label{eq:poissonL}
\end{equation}
where $\hat K_a$ are Kraus operators for the whole channel and
$\sum_a\hat K_a^\dagger\hat K_a=\Id$.
Equivalently, subtract the input state from the probabilistic
update and divide by the elapsed time:
\begin{equation}
 \frac{\hat\rho(t+dt)-\hat\rho(t)}{dt}
       =\gamma\bigl[\E(\hat\rho(t))-\hat\rho(t)\bigr].
 \label{eq:time_step_derivative}
\end{equation}
The limit of repeated small steps is the continuous-time
semigroup generated by this right-hand side. The Kraus expression
in Eq.~\eqref{eq:poissonL} independently verifies that this is a
valid Lindblad generator, not just a formal differential equation.
Its semigroup has the
explicit completely positive representation
\begin{equation}
 e^{t\mathcal L_{\rm cyc}}
   =e^{-\gamma t}\sum_{n=0}^{\infty}
            \frac{(\gamma t)^n}{n!}\E^n. \label{eq:poisson}
\end{equation}
It applies a Poisson-distributed number of complete cycles and has
the same stationary space as $\E$. Physically, the deterministic
sequence inside a cycle is unchanged; only the times at which
complete cycles occur are randomized. A mode with multiplier $\mu$
therefore acquires the averaged factor
\begin{equation}
 e^{-\gamma t}\sum_{n=0}^{\infty}
       \frac{(\gamma t)^n}{n!}\mu^n
       =e^{\gamma t(\mu-1)}.
 \label{eq:poisson_multiplier}
\end{equation}
This also explains why phase oscillations can become damping:
averaging over different numbers of steps can cancel them.
The eigenvalues and gap obey
\begin{align}
 \lambda&=\gamma(\mu-1),\nonumber\\
 \delta_{\mathcal L_{\rm cyc},\mathcal V}
   &=\gamma\min_{\mathcal V_0}(1-\operatorname{Re}\mu)
     \ge\gamma g. \label{eq:poisson_bound}
\end{align}
The converse fails: Eq.~\eqref{eq:flip} has $g=0$ but a positive
Poissonized decay rate. More generally, eigenvalues elsewhere on
the unit circle have negative real parts after subtracting 1.
For the lazy flip channel the nonstationary eigenvalues are
$1-2p,1-p,1-p$; it mixes for $0<p<1$, although the original
$p=1$ channel does not.

For the reset covariance spectrum, the dominant multiplier $r^2$
is real, nonnegative, and present in the spectrum. It saturates
both inequalities: its probabilistic multiplier is
$1-p+pr^2$, and its Poissonized exponent is $\gamma(r^2-1)$.
All other covariance multipliers obey the previously established
bounds, so none can be slower. This gives the exact identities
\begin{equation}
 g_C(p)=p\,g_C,\qquad
 \delta_{\mathcal L_{\rm cyc},C}=\gamma g_C. \label{eq:exact_poisson}
\end{equation}
By App.~\ref{sec:spectrum}, the same equalities hold in the
neutral and even many-body sectors. The continuous-time rate in
Eq.~\eqref{eq:exact_poisson} is $\gamma g_C$, not $\gamma\kappa_C$.
A per-step gap proportional to $dt$ tends to zero when the clock
step shrinks without implying slow relaxation in physical time.
For example, applying a channel with probability $\gamma dt$
per tick means that most sufficiently short ticks do nothing.
The number of opportunities per unit time grows as $1/dt$ and
compensates the small gap per tick. Poissonization measures this
physical-time decay; $\kappa_C$ instead describes exactly one
completed reset cycle per discrete step.

\subsection{Finite-time semigroups and limits of the correspondence}
If a relaxing Lindbladian is given first, its finite-time channel
$\E_\tau=e^{\tau\mathcal L}$ has multipliers $e^{\tau\lambda}$.
In the same invariant sector,
\begin{equation}
 g(\E_\tau)=1-e^{-\tau\delta_{\mathcal L}},\qquad
 \kappa(\E_\tau)=\tau\delta_{\mathcal L}. \label{eq:semigroup}
\end{equation}
This spectral mapping statement differs from Poissonizing an
arbitrary channel. In general $\E$ need not admit a Lindbladian
logarithm; semigroup-embedding restrictions already occur for
qubit channels~\cite{PRZ}. Equation~\eqref{eq:poissonL} always
constructs a generator, but its exponential is not generally $\E$.

Three qualifications are important. First, whole-cycle Kraus
operators can extend across the system. A local gate sequence does
not automatically make $\mathcal L_{\rm cyc}$ a finite-range
Lindbladian, especially when the sequential depth grows with size.
Weakening each elementary reset separately instead produces, to
first order, $\sum_j\gamma_j(\E_j-\mathrm{Id})$. This generator
need not share the stationary state of the ordered full-reset
product. They do share a state if every reset fixes it individually,
but such a common fixed state is an additional assumption.

Second, finite-dimensional eigenvalue identities do not bound
transient amplification or all mixing prefactors. The projector
product and its induced maps can be nonnormal, and defective
blocks introduce powers of time multiplying exponentials.
Third, thermodynamic gap preservation requires a size-independent
lower bound on $g$, with $p$ bounded away from zero or $\gamma$
fixed. Positivity at each finite size is insufficient.
Local mixing and clustering theorems such as those of Kastoryano
and Eisert~\cite{KE} have further regularity and reversibility or
stronger mixing assumptions. They cannot be imported for the
whole-cycle construction without checking those hypotheses.

\section{Trajectory dynamics and relation to the literature}
\label{sec:trajectory}
\subsection{Conditioning versus discarding the record}
For a complete measurement record $\bm m$ let $\hat V_{\bm m}$ be
the corresponding Kraus product. Then
\begin{align}
 p_{\bm m}&=\Tr(\hat V_{\bm m}\hat\rho_0\hat V_{\bm m}^\dagger),
 &\hat\rho_{\bm m}&=\frac{\hat V_{\bm m}\hat\rho_0
                   \hat V_{\bm m}^\dagger}{p_{\bm m}},\nonumber\\
 \overline{\hat\rho}_n&=\sum_{\bm m}p_{\bm m}\hat\rho_{\bm m}
                     =\E^n(\hat\rho_0).&& \label{eq:records}
\end{align}
Linear observables obey
$\sum_{\bm m}p_{\bm m}\Tr(\hat\rho_{\bm m}\hat O)
=\Tr(\overline{\hat\rho}_n\hat O)$.
Entropies do not: generally
\begin{equation}
 \sum_{\bm m}p_{\bm m}S((\hat\rho_{\bm m})_R)
       \ne S((\overline{\hat\rho}_n)_R), \label{eq:entropy}
\end{equation}
where $R$ is a spatial region. Different unravelings can yield the
same averaged channel but different conditioned entanglement.
The failure of entropy to commute with averaging is elementary.
An equal ensemble of an empty and an occupied one-mode state
has zero entropy in each member, but its averaged state
$\Id/2$ has entropy $\log 2$. This is a statement about ensemble
mixing, not an entanglement claim for that one mode; the same
nonlinearity prevents replacing averaged subsystem entropies by
the entropy of the averaged subsystem state.
Accordingly, algebraic relaxation of trajectory-resolved
entanglement is not, by itself, evidence that the one-copy
outcome-averaged channel has zero gap. A relation would require an
additional argument connecting that nonlinear record-dependent
quantity to the relevant channel or a replicated evolution.

This distinction separates the conditioned purification discussed
in BPJ Fig.~9 from the averaged correlation matrix and occupation
spectrum in BPJ Fig.~7~\cite{BPJ}. Spatial algebraic correlations,
temporal decay of averaged linear observables, and temporal decay
of conditioned entanglement are three different diagnostics.
Even a temporal power law over a finite window need not determine
the asymptotic finite-size rate if transients or crossovers remain.
None of these distinctions negates a measured power law; they
identify which spectral claim that measurement can support.

\subsection{Local dissipative examples and scope of existing results}
A relevant finite-range example is Bardyn \emph{et al.}~\cite{Bardyn},
Eqs.~(68)--(69) and Fig.~6. Their competing one-dimensional jump
families are
\begin{equation}
 \hat L_n^{(\ell)}=\tfrac12\bigl[
   (\hat a_n^\dagger+\hat a_{n+\ell}^\dagger)
   +(\hat a_n-\hat a_{n+\ell})\bigr],\qquad \ell=1,2.
 \label{eq:bardyn}
\end{equation}
With their normalized relative-rate convention, the bulk damping
spectrum remains flat and gapped as the bulk purity gap closes at
equal competing strengths. This is a finite-range real-space
construction, not merely independent control of nonlocal momentum
modes. The bulk qualifier is essential: open boundaries can carry
zero-damping Majorana edge modes. The example therefore does not
establish a gapped full open-boundary generator. It establishes
that locality alone does not equate the two bulk gaps.

Relations can reappear with extra structure. In the Gaussian
setting, damping and noise jointly determine the stationary
covariance; special dark-state or frustration-free conditions can
lock these quantities~\cite{Mittal}. Such restrictions are not a
general identity between stationary purity and relaxation gaps.
Likewise, the symmetry-class correspondence between quadratic
Lindbladians and their unique Gaussian stationary states~\cite{Mao}
is a statement about symmetry constraints, not equality of gaps.
Its linear-jump assumptions do not automatically cover the
number-dephasing term in Eq.~\eqref{eq:idplusL}.

Discrete-channel studies also treat dynamical and invariant-state
spectra as separate diagnostics. Spectral-support transitions in
random Kraus maps~\cite{Sa} should not be conflated with a relaxation
gap closing or a fermionic purity-gap crossing. Noisy Floquet
channels provide examples with maximally mixed stationary states
and nontrivial relaxation, while emphasizing the order of weak-noise
and thermodynamic limits~\cite{YS}. Broader random-channel results
give further context for spectra and invariant states~\cite{Kukulski},
not a universal relation between their gaps. For the present reset
model, the precise many-body statement remains the hierarchy proof
in App.~\ref{sec:spectrum}.

\section{Numerical illustration and verification}
\label{sec:numerics}
\subsection{Common protocol and postprocessing}
We use two completed CPU campaigns of the canonical
\texttt{run\_markov\_channel} implementation. Both use
$\alpha_1=1,3$, $\alpha_2=30$, $n_{\rm shell}=1$, hard-wall
support truncation, all slabs active, periodic boundaries, zero
twist, $X$ trial orbitals, perfect correction, and double-precision
complex arithmetic. Within each cell the order is
$A+,A-,B+,B-$, with the raster-$y$ schedule ($x$ outer, $y$ inner).
The initial correlation matrix is $\id/2$. These are exact
outcome-averaged covariance calculations, not trajectory samples.

Each run contains 61 cycles, with saved occupation spectra at
$0,1,5,10,20,40,60,61$. We report cycle 61. For all 24 runs,
the last five successive Frobenius changes are below $10^{-10}$,
and the full and twirled purity gaps change by less than
$2\times10^{-10}$ between cycles 60 and 61. This is a numerical
stationarity check, not a rigorous norm bound on the error to the
fixed point of a nonnormal map.

The translation-twirled matrix is the postprocessed average
\begin{equation}
 \mathcal T_y(G)=\frac1{N_y}\sum_{r=0}^{N_y-1}
        T_y^rG(T_y^\dagger)^r. \label{eq:twirl}
\end{equation}
It is block diagonal in $k_y=2\pi m/N_y$; its spectrum is the union
of the spectra of those momentum blocks, not an average of the
original eigenvalues. To see this, in a translation eigenbasis
each matrix element between $k_y$ and $k_y'$ is multiplied by
$N_y^{-1}\sum_r e^{ir(k_y-k_y')}=\delta_{k_y,k_y'}$.
Twirling removes couplings between distinct momenta but retains
the matrix within each momentum block; those retained matrices
must then be diagonalized. Twirling does not alter the simulated
schedule or produce another dynamics run. In particular,
$\mathcal T_y(G_\steady)$ need not itself be a fixed point of the
raster-ordered channel.

\subsection{Fixed-width systems}
The fixed-width campaign uses $N_x=20$ and
$N_y=20,40,60,80,100$, with inclusive walls at $x=5,15$.
Table~\ref{tab:fixed} reports the full and twirled purity gaps.
The $\alpha_1=1$ values are small and nonmonotone; the
$\alpha_1=3$ values remain near $0.986$. Occupation half-distances
$\Delta$ are one half of every tabulated entry.

\begin{table}[tbp]
\caption{Fixed width $N_x=20$: stationary purity gaps
$\Delta_{\rm pur}=\min|1-2n|$. ``Twirled'' refers to
Eq.~\eqref{eq:twirl}. All entries are regenerated from the saved
cycle-61 occupation spectra, with six decimal places.}
\label{tab:fixed}
\begin{ruledtabular}
\begin{tabular}{rcccc}
 &\multicolumn{2}{c}{$\alpha_1=1$}
 &\multicolumn{2}{c}{$\alpha_1=3$}\\
 $N_y$ & Full & Twirled & Full & Twirled\\
 % Generated from verified cycle-61 occupation arrays by build_note.py.
20 & 0.086790 & 0.092435 & 0.985595 & 0.985567 \\
40 & 0.004357 & 0.003950 & 0.985659 & 0.985703 \\
60 & 0.021813 & 0.022742 & 0.985667 & 0.985748 \\
80 & 0.005699 & 0.005474 & 0.985668 & 0.985771 \\
100 & 0.009523 & 0.009948 & 0.985668 & 0.985784 \\

\end{tabular}
\end{ruledtabular}
\end{table}

Separate matched spectral calculations exist for these ten cases.
Their covariance channel gaps decrease from $0.851086$ to
$0.835675$ for $\alpha_1=1$, and from $0.961228$ to $0.957611$
for $\alpha_1=3$, between $N_y=20$ and 100. These are $g_C$,
not the logarithmic rates $\kappa_C$. Thus small stationary purity
gaps coexist with large finite-size relaxation gaps in this
fixed-width sequence. Increasing $N_y$ at fixed $N_x$ is not the
two-dimensional fixed-aspect-ratio thermodynamic limit.

For a concrete reading of the numbers, at
$(N_x,N_y,\alpha_1)=(20,20,1)$ the full purity gap $0.086790$
means that the closest occupation is a distance $0.043395$
from half filling. Independently, $g_C=0.851086$ means that the
slowest covariance multiplier is approximately $0.148914$
per cycle. The first number locates an occupation in the final
state; the second gives the asymptotic decay factor of a perturbation.
Their sizes need not track one another.

The twirled spectra provide evidence of momentum-grid effects:
the relevant branch at $\alpha_1=1$ comes closest to half filling
near $|k_y|\simeq0.05\pi$. Grids with $N_y=40$ and 80 sample that
region especially closely. The overlapping branch shapes and
nonmonotone minima are consistent with commensurability, but do
not prove an exact crossing between grid points or a thermodynamic
purity-gap closing. No interpolation or fit is used here.

\subsection{Square systems and provenance}
The square campaign uses $N_x=N_y=L=20,30,40,50,60,70,80$.
The inclusive wall positions are $\lfloor L/4\rfloor$ and
$\lfloor3L/4\rfloor$; their integer rounding is retained.
Table~\ref{tab:square} again shows small, nonmonotone
$\alpha_1=1$ purity gaps and large $\alpha_1=3$ gaps. These
data vary both dimensions and must not be pooled with the
fixed-width sequence as though the geometries were identical.
This new square \emph{dynamics} sweep did not compute channel
gaps. Table~\ref{tab:square} makes no claim about their size dependence.

\begin{table}[tbp]
\caption{Square systems $N_x=N_y=L$: full and postprocessed
twirled purity gaps, with the same definitions and precision as
Table~\ref{tab:fixed}. These are stationary occupation results,
not a square-system channel-gap sweep.}
\label{tab:square}
\begin{ruledtabular}
\begin{tabular}{rcccc}
 &\multicolumn{2}{c}{$\alpha_1=1$}
 &\multicolumn{2}{c}{$\alpha_1=3$}\\
 $L$ & Full & Twirled & Full & Twirled\\
 % Generated from verified cycle-61 occupation arrays by build_note.py.
20 & 0.086790 & 0.092435 & 0.985595 & 0.985567 \\
30 & 0.024280 & 0.026040 & 0.985247 & 0.985250 \\
40 & 0.004357 & 0.003950 & 0.985086 & 0.985106 \\
50 & 0.022346 & 0.022455 & 0.984999 & 0.985031 \\
60 & 0.021813 & 0.022741 & 0.984946 & 0.984988 \\
70 & 0.005190 & 0.005699 & 0.984912 & 0.984961 \\
80 & 0.005699 & 0.005474 & 0.984889 & 0.984943 \\

\end{tabular}
\end{ruledtabular}
\end{table}

The accompanying build script verifies byte counts and SHA-256
digests for all 24 result/completion pairs, and independently
recomputes both purity-gap columns from saved occupation spectra.
It also verifies the ten matched fixed-width spectral pairs.
The data roots are the October 5, 2026 runs
\texttt{steady\_purity\_gap\_v1} (18:41:36 UTC) and
\texttt{square\_steady\_purity\_gap\_v1} (22:54:46 UTC).
Exact paths, configuration and source hashes, convergence checks,
and unrounded values are retained in the machine-readable
validation record. Small-Fock-space checks test the reset identity,
covariance convention, exterior-power spectra, operator sectors,
and continuous-time examples. No new dynamics, fits, or
thermodynamic extrapolations enter this consolidation.

\begin{thebibliography}{13}
\bibitem{BPJ} A. Bhuiyan, H. Pan, and C.-M. Jian,
\href{https://doi.org/10.1103/nk38-ygyq}{Free-fermion dynamics with measurements:
Topological classification and adaptive preparation of topological states},
Phys. Rev. Research \textbf{8}, 023147 (2026);
\href{https://arxiv.org/abs/2507.13437}{arXiv:2507.13437}.
\bibitem{Burgarth} D. Burgarth, G. Chiribella, V. Giovannetti,
P. Perinotti, and K. Yuasa,
\href{https://doi.org/10.1088/1367-2630/15/7/073045}{Ergodic and mixing
quantum channels in finite dimensions}, New J. Phys. \textbf{15}, 073045 (2013);
\href{https://arxiv.org/abs/1210.5625}{arXiv:1210.5625}.
\bibitem{Wolf} M. M. Wolf,
\href{https://mediatum.ub.tum.de/doc/1701036/document.pdf}{Quantum Channels
\& Operations: Guided Tour}, lecture notes (2012), Sec.~7.1.
\bibitem{BZ} T. Barthel and Y. Zhang,
\href{https://doi.org/10.1088/1742-5468/ac8e5c}{Solving quasi-free and quadratic
Lindblad master equations for open fermionic and bosonic systems},
J. Stat. Mech. (2022) 113101;
\href{https://arxiv.org/abs/2112.08344}{arXiv:2112.08344}.
\bibitem{DFP} B. Dierckx, M. Fannes, and M. Pogorzelska,
\href{https://doi.org/10.1063/1.2841326}{Fermionic quasi-free states and maps
in information theory}, J. Math. Phys. \textbf{49}, 032109 (2008);
\href{https://arxiv.org/abs/0709.1061}{arXiv:0709.1061}.
\bibitem{PRZ} Z. Pucha\l{}a, \L{}. Rudnicki, and K. \.{Z}yczkowski,
\href{https://doi.org/10.1016/j.physleta.2019.04.057}{Pauli semigroups and
unistochastic quantum channels}, Phys. Lett. A (2019);
\href{https://arxiv.org/abs/1903.12448}{arXiv:1903.12448}.
\bibitem{KE} M. J. Kastoryano and J. Eisert,
\href{https://doi.org/10.1063/1.4822481}{Rapid mixing implies exponential
decay of correlations}, J. Math. Phys. \textbf{54}, 102201 (2013);
\href{https://arxiv.org/abs/1303.6304}{arXiv:1303.6304}.
\bibitem{Bardyn} C.-E. Bardyn, M. A. Baranov, C. V. Kraus, E. Rico,
A. Imamoglu, P. Zoller, and S. Diehl,
\href{https://doi.org/10.1088/1367-2630/15/8/085001}{Topology by dissipation},
New J. Phys. \textbf{15}, 085001 (2013);
\href{https://arxiv.org/abs/1302.5135}{arXiv:1302.5135}.
\bibitem{Mittal} R. Mittal, T. Zander, J. Lang, and S. Diehl,
\href{https://doi.org/10.1103/nj2d-l6pz}{Fermion quantum criticality far
from equilibrium}, Phys. Rev. X \textbf{16}, 011069 (2026);
\href{https://arxiv.org/abs/2507.14318}{arXiv:2507.14318}.
\bibitem{Mao} L. Mao and F. Yang,
\href{https://doi.org/10.1103/p8xr-4yzl}{Symmetry classification
correspondence between quadratic Lindbladians and their steady states},
Phys. Rev. B \textbf{112}, 085136 (2025);
\href{https://arxiv.org/abs/2503.10763}{arXiv:2503.10763}.
\bibitem{Sa} L. S\'a, P. Ribeiro, T. Can, and T. Prosen,
\href{https://doi.org/10.1103/PhysRevB.102.134310}{Spectral transitions and
universal steady states in random Kraus maps and circuits},
Phys. Rev. B \textbf{102}, 134310 (2020);
\href{https://arxiv.org/abs/2007.04326}{arXiv:2007.04326}.
\bibitem{YS} T. Yoshimura and L. S\'a,
\href{https://www.nature.com/articles/s41467-024-54164-7}{Robustness of
quantum chaos and anomalous relaxation in open quantum circuits},
Nat. Commun. \textbf{15}, 9808 (2024);
\href{https://arxiv.org/abs/2312.00649}{arXiv:2312.00649}.
\bibitem{Kukulski} R. Kukulski, I. Nechita, \L{}. Pawela, Z. Pucha\l{}a,
and K. \.{Z}yczkowski,
\href{https://doi.org/10.1063/5.0038838}{Generating random quantum channels},
J. Math. Phys. \textbf{62}, 062201 (2021);
\href{https://arxiv.org/abs/2011.02994}{arXiv:2011.02994}.
\end{thebibliography}
\end{document}
